\section{模型假设}

结合题面条件与数据的实际结构，本文作如下假设。凡属题面未规定而由本文约定的做法，均在条目中直接说明，不以隐含默认的方式使用。

\begin{enumerate}[leftmargin=2.4em,itemsep=2pt]
  \item \textbf{时间离散与对齐}：电价、负载与光伏功率均以 10 分钟为时段、每日 144 段；附件时间戳为所属时段的\emph{左端点}，第 $i$ 个数据点代表区间 $[10i,10i+10)$ 分钟，与结果模板的 144 个时段按位置一一对应。据此计划窗为 0:10 至 24:10，模板标称的“0:00--4:00”等标签对应物理区间整体前移 10 分钟，属模板固有相位，全文表格与图注均按此约定统一。
  \item \textbf{时段内恒定}：每个时段内电价与功率视为常数，时段电量等于功率乘以 $\Delta t=1/6$ h；反之，以电量表示的决策变量本身已含时间因子，费用表达式中不再乘以 $\Delta t$。
  \item \textbf{储能模型}：储能按常效率模型描述，充、放电各按 90\% 计（往返效率 81\%），静置时不发生自放电；容量运行区间为 $[1200,10800]$ kWh，初始储电量为 6000 kWh。
  \item \textbf{功率上限的作用位置}：题面“最大充放电功率 5000 kW”同时作用于并网点侧与电池侧、取更严者，故充电量上界为 $5000\Delta t=833.3333$ kWh、放电量上界为 $5000\Delta t\,\eta=750.0000$ kWh。该项为题面参数的解释约定，本文不将其表述为文献结论。
  \item \textbf{弃光变量保留}：光伏盈余中未被负载与储能消纳的部分记为弃光量 $s$，其目标系数为零（无上网售电收益），但变量必须保留，以便在盈余超出储能吸收能力的时段闭合功率平衡。
  \item \textbf{购电无功率上限、无售电收益}：外网购电量非负且不设上界；光伏余电不上网，储能也不向电网售电。当分析中引入购电上限或售电价格时，均视为独立的结构性对照情景单独报告。
  \item \textbf{费用范围}：目标只计购电费用，不引入折旧、循环寿命或退化成本（题面未给系数，无法标定）。紧急购电按交易时刻电价的 5 倍计价；问题三与链 \emph{4-3} 中，计划高于最终的部分按 0.5 倍价、最终高于计划的部分按 1.5 倍价计价，两段均为成本项而非退款，且以当天 0:00 计划为唯一基准、单次结算、不回溯。
  \item \textbf{信息非预期性}：任一决策时刻只能使用该时刻已发布或已实现的数据。问题三与链 \emph{4-3} 中，0:00 计划层只用 0:00 预报，6:00、12:00、18:00 调整层只用该时刻已发布预报加已实现状态；实际光伏只进入事后结算层，不影响任何决策量与储能轨迹。
  \item \textbf{价格可观测性}：在附件 4 的波动电价情景下，0:00 制定当天计划时当天实时电价尚未发生、不可观测，决策层只能使用由历史价格构造的预测价，而费用结算一律使用实际电价；由此产生的决策—结算差异单独量化为“价格预报误差的代价”。
  \item \textbf{供电可靠性由等式平衡保证}：每个时段采用等式功率平衡并显式包含紧急购电与弃光变量，由此“微网提供的电能不低于小区负载”自动成立、充电不挤占负载；模型不进一步建模切负荷。
  \item \textbf{数据可用性与预报边界}：附件 2--4 在 2025 全年无缺日、无重复，附件 3 的日期列存在前向填充结构，读取时按前向填充还原；题面只提供未来 24 小时的整点光伏预报，本文不假设可获得更长时间跨度或更高时间分辨率的预报。
\end{enumerate}

此外，由于最优调度面可能不唯一（储电量上下限两端均可能活跃），本文统一约定一条\textbf{解的选择规则}：所有模型、所有层先在主目标最优面上取解，再在其中最小化储能吞吐量（字典序规则）。验收一律以约束残差、结构恒等式、目标值与交付表格为准，不比对逐点解的唯一性。
